\section{Quick Start}\label{sec:quick-start}

\begin{enumerate}
  \item Connect an oscillator or another signal to the input. The illuminated
  \textbf{Run} button indicates that Spectre is acquiring new spectra.
  Keep your existing connection to the mixer or audio output if you want to
hear it.
  \item A new module uses $+6\,\mathrm{dB}$ input gain, a Flattop window,
  and no time or frequency smoothing. For comparisons with Fourier, set
  both modules to the same input gain; Fourier starts at $0\,\mathrm{dB}$.
  \item A steady sine wave produces a horizontal band. Harmonics appear as
  additional bands above it; short broadband sounds produce vertical marks.
  \item Drag the two limits on the vertical \textbf{Color} bar to reveal quiet detail or
  distinguish bright regions, including while frozen. Set \textbf{Slope} to zero when you
  want to compare magnitudes without frequency-dependent display weighting.
  The default is $+4.5\,\mathrm{dB/oct}$.
  \item Press \textbf{Run} to freeze the history, then hover over the plot to
  inspect a frequency and its magnitude. Press it again to resume.
\end{enumerate}

\paragraph{Reading levels} Spectre displays spectral magnitudes, not total
signal RMS, loudness, or sound pressure level. Window choice, frequency
alignment with FFT bins, and smoothing affect the height and width of a peak.
The color map has a finite range; reaching its end color does not mean Spectre
has clipped your audio path. The \textbf{Raw} hover readout measures stored magnitudes before Slope
weighting; \textbf{Color} measures the weighted, interpolated image sample.

\paragraph{Choosing settings} Use little or no \textbf{Average} to follow
changes more closely. Increase \textbf{Average} for a steadier image and
\textbf{Smooth} for a broader view of tonal balance. FFT length and hop size
are fixed; cropping the frequency range enlarges existing detail without
increasing resolution. The factory presets provide examples of these choices.
